\section{凸松弛把 NP-hard 问题变成可解问题}
\label{sec:09-Convex-Optimization/09-Complexity-and-Approximation/02}

选址模型跑了一夜没有结束。模型里有一行 \(x\in\set{0,1}^{n}\)，意思是每座候选仓库要么建要么不建。你把这一行改成 \(0\le x\le1\)，重新提交，四秒返回：目标值 \(187.3\)，其中三个仓库拿到的值是 \(0.4\)、\(0.6\)、\(0.5\)。

这个方案没法执行，没有人会去建半座仓库。但屏幕上那个数字不是废品。它是一个下界：真实的最优成本不可能低于 \(187.3\)。手上任何一个可执行的方案，比如某个贪心排出来的 \(212.0\)，现在都带上了一句可核对的话——距最优不超过 \(13\%\)。接下来的全部工作，是把那几个小数变回 0 和 1，以及说清楚这一步会损失多少。

\subsection{放宽约束，最优值只会变得更好看}

把可行集换成一个包含它的更大的集合，这个动作叫松弛。原问题

\[
\begin{aligned}
\text{minimize}\quad&c^{\trans}x\\
\text{subject to}\quad&Ax\le b,\\
&x\in\set{0,1}^{n}
\end{aligned}
\]

的最后一行换成 \(0\le x\le1\)，可行集从有限个顶点扩成一个多面体，问题从 NP-hard 的整数规划变成多项式时间可解的线性规划。

界的方向是自动的，不需要任何额外条件：可行集变大，最小值只会变小或不变，所以

\[
p^\star_{\text{松弛}}\le p^\star.
\]

这句话的分量与\hyperref[sec:09-Convex-Optimization/05-Duality/02]{``对偶间隙何时闭合''}里的弱对偶是同一类——两者都无条件成立，都给下界，实际上 Lagrange 对偶可以看成松弛的另一种写法。至于这个下界紧不紧，得另外论证，定理不管这一层。

有一种情形是白送的：松弛解本身恰好满足被松掉的约束。跑出来的 \(x\) 全是 0 和 1，那它在原可行集里，又达到了一个下界，于是它就是原问题的最优解，并且附赠一份最优性证明。这是唯一一种能免费拿到``不可能更好''的场景，值得每次都检查一遍。

\subsection{间隙的来源是你写下的那个多面体，不是凸包}

\begin{center}
\begin{tikzpicture}[x=1.05cm,y=1.05cm,font=\small,>=stealth]
  \begin{scope}
    \fill[black!8] (0,0) -- (3.6,0) -- (3.6,1.2) -- (1.2,3.6) -- (0,2.2) -- cycle;
    \draw[dashed,black!55] (0,0) -- (3.6,0) -- (3.6,1.2) -- (1.2,3.6) -- (0,2.2) -- cycle;
    \fill[blue!12] (0,0) -- (3,0) -- (3,1) -- (1,3) -- (0,2) -- cycle;
    \draw[thick,blue!55!black] (0,0) -- (3,0) -- (3,1) -- (1,3) -- (0,2) -- cycle;
    \foreach \x/\y in {0/0,1/0,2/0,3/0,0/1,1/1,2/1,3/1,0/2,1/2,2/2,1/3}
      \fill[black!65] (\x,\y) circle (1.7pt);
    \draw[dashed,black!45] (3.9,-0.5) -- (1.8,3.0);
    \draw[dashed,black!45] (4.62,-0.5) -- (2.52,3.0);
    \fill[red!70!black] (3,1) circle (2.6pt);
    \fill[red!70!black] (3.6,1.2) circle (2.6pt);
    \node[font=\footnotesize,red!60!black,left] at (2.88,0.95) {\(x^\star\)};
    \node[font=\footnotesize,red!60!black,right] at (3.7,1.35) {\(x_{\text{松弛}}\)};
    \draw[<->,black!75] (3.9,-0.5) -- (4.62,-0.5);
    \node[font=\footnotesize,black!75,right] at (4.68,-0.5) {间隙};
    \draw[->,black!70] (0.35,2.95) -- (1.1,3.4);
    \node[font=\footnotesize,black!70,above left] at (1.1,3.36) {\(-c\)};
    \node[font=\footnotesize] at (2.0,-0.95) {写下来的松弛比凸包大};
  \end{scope}
  \begin{scope}[shift={(6.4,0)}]
    \fill[blue!12] (0,0) -- (3,0) -- (3,1) -- (1,3) -- (0,2) -- cycle;
    \draw[thick,blue!55!black] (0,0) -- (3,0) -- (3,1) -- (1,3) -- (0,2) -- cycle;
    \draw[dashed,black!55] (0,0) -- (3,0) -- (3,1) -- (1,3) -- (0,2) -- cycle;
    \foreach \x/\y in {0/0,1/0,2/0,3/0,0/1,1/1,2/1,3/1,0/2,1/2,2/2,1/3}
      \fill[black!65] (\x,\y) circle (1.7pt);
    \draw[dashed,black!45] (3.6,0) -- (1.8,3.0);
    \fill[red!70!black] (3,1) circle (2.6pt);
    \node[font=\footnotesize,red!60!black,right] at (3.12,1.15) {\(x^\star=x_{\text{松弛}}\)};
    \node[font=\footnotesize] at (2.0,-0.95) {两者重合，间隙为零};
  \end{scope}
\end{tikzpicture}

\emph{黑点是原问题的可行解，蓝色区域是它们的凸包。左边灰色的那圈是你在模型里实际写下的松弛，它包住凸包但更大，于是最优点落在一个分数顶点上，两条平行的等值线之间那段就是松弛间隙。右边这两个多面体重合，线性规划的最优顶点自动是一个整点。}
\end{center}

图里有一件事值得停一下：间隙并不来自``把离散换成连续''这个动作本身。在凸包上做线性优化，最优值一定在某个顶点取到，而凸包的顶点全是原来的可行点，所以在凸包上求解与在离散集上求解给出同一个答案（\hyperref[sec:09-Convex-Optimization/02-Convex-Sets/01]{``线段、凸组合与凸包''}里已经有过这个观察）。间隙来自另一件事：凸包通常需要指数多条不等式才描述得出来，你写不下它，只能写一个更宽松的多面体，比如那个盒子与 \(Ax\le b\) 的交。间隙的宽窄，量的是你写下的那圈边界离凸包有多远。

于是零间隙的条件也清楚了：写下的多面体的每个顶点都是整点。线性规划的顶点是若干条约束的交点，由 \(A\) 的某个子矩阵解出来，若 \(A\) 的所有子式都取 \(0\) 或 \(\pm1\)（这样的矩阵称为全幺模），而 \(b\) 是整向量，那么每个顶点的坐标都是整数，松弛白送最优解。指派问题、二分图匹配、网络流的约束矩阵恰好都是全幺模的——上一篇开头那个可以答应``最优''的派单需求，底层就是这条性质在起作用。它不是运气，是模型结构的一项可以查验的属性。

\subsection{把小数变回可行解，有三种做法}

拿到分数解之后的这一步是纯工程活，做法的差别直接决定你能对外说什么。

最直接的是阈值取整：大于某个门槛的置 1，其余置 0。风险在于取整之后可能不可行。顶点覆盖是个正面例子：LP 松弛要求每条边的两个端点满足 \(x_i+x_j\ge1\)，于是至少有一个端点的值不小于 \(1/2\)，把所有不小于 \(1/2\) 的顶点都选上，每条边必被覆盖。这个舍入还带一个比例保证：舍入只把每个 \(x_i\) 放大不超过两倍，所以解的规模不超过 LP 值的两倍，而 LP 值本身不超过最优值，两步串起来就是 2 倍近似。条件对账在于那句``至少一个端点不小于 \(1/2\)''——它用到了约束的具体形状，换一个问题就不成立了。容量约束下的四舍五入常常直接把总量顶破，那时取整给出的东西连方案都算不上。

第二种是随机舍入：把 \(x_i\in[0,1]\) 读成概率，独立地以 \(x_i\) 的概率把第 \(i\) 项置 1。期望值这一步是精确的，\(\E[\text{目标}]\) 正好等于松弛的目标值，剩下的分析是控制可行性与偏差。最大割的 Goemans--Williamson 算法是这套路子的样板：先把 \(x_i\in\set{-1,1}\) 松弛成单位向量 \(v_i\)，解一个半正定规划（见\hyperref[sec:09-Convex-Optimization/04-Convex-Problems/04]{``半正定规划优化矩阵不等式''}），再取一个随机方向 \(r\)，令 \(x_i=\sign(\ip{v_i}{r})\)。一条边被切开的概率是它两端向量夹角除以 \(\pi\)，与松弛里那条边的贡献逐项比较，比值的最小值

\[
\alpha=\min_{0\le\theta\le\pi}\frac{2}{\pi}\cdot\frac{\theta}{1-\cos\theta}\approx0.878
\]

就是期望近似比。整段论证的重量都压在``逐边比较''这一步上，它把一个全局的期望拆成了每条边各自的保证。

第三种是不舍入，把界拿去剪枝。分支定界在某个分数变量上劈成 \(x_i=0\) 与 \(x_i=1\) 两支，各自重新求松弛。某一支的下界若已经超过当前最好的可行方案，整支连同它下面所有分叉一起丢掉。这条路最终能给出精确最优解，换回来的是最坏情形的指数时间，而实际能不能跑完，几乎完全取决于松弛紧不紧。求解器日志里那个缓慢下降的下界，就是这些松弛在一层层收紧。

\subsection{\(\ell_0\) 到 \(\ell_1\) 是这套想法最常用的一次}

风控模型要求上线时最多用 30 个特征，采集成本与可解释性都卡在这里。写成优化问题就是

\[
\begin{aligned}
\text{minimize}\quad&\norm{Xw-y}_2^2\\
\text{subject to}\quad&\norm{w}_0\le30,
\end{aligned}
\]

其中 \(\norm{w}_0\) 数的是非零分量的个数。它不是范数，也不连续，这条约束把问题变成了在 \(\binom{800}{30}\) 个支撑集上的搜索——子集选择是 NP-hard 的。

把``非零个数最少''换成``绝对值之和最小''，即用 \(\norm{w}_1\) 顶替 \(\norm{w}_0\)，问题落进凸的范畴，几秒钟解完。这个替换不是随手挑的：在盒子 \(\norm{w}_\infty\le1\) 上，\(\norm{w}_1\) 恰是 \(\norm{w}_0\) 的凸包络，也就是所有落在 \(\ell_0\) 下方的凸函数里最大的那一个。凸包络这个词说的正是上一小节图里那件事，只不过从可行集换成了函数图像。为什么这个替换还顺带产生稀疏解，而不是仅仅把系数压小，\hyperref[sec:09-Convex-Optimization/07-Applications/02]{``\(L_1\) 正则为何产生稀疏''}用次微分给过解释。

同一手法在矩阵上有个孪生兄弟：\(\rank(X)\) 换成核范数 \(\norm{X}_{*}=\sum_i\sigma_i(X)\)，秩的最小化变成半正定规划，这是\hyperref[sec:13-Machine-Learning/04-Dimensionality-Reduction/03]{``矩阵补全从缺失条目到低秩正则''}的出发点。奇异值向量的 \(\ell_1\) 范数替代它的 \(\ell_0\) 范数，一字不差的同一句话。

值得注意的是这里的``变回可行解''那一步长什么样。\(\ell_1\) 解通常给出一批很小但非零的系数，工程上的做法是按阈值截断取支撑集，再在这个支撑集上重新做一次无正则的最小二乘。第二次拟合是必要的：\(\ell_1\) 罚项在压零的同时也把留下来的系数往零方向拖了一截，不去掉这个收缩，预测值会系统性偏低。松弛给的是支撑集这条信息，不是系数本身。

\subsection{松弛的紧度取决于模型怎么写，而它有天花板}

同一个组合问题有很多种等价的整数规划写法，它们的 LP 松弛却可以差出很远。用大 \(M\) 系数把逻辑条件写成线性约束时，\(M\) 取得越大松弛越松，分支定界越跑不动；把一条聚合约束拆成若干条更强的局部约束，可行的整数点一个不少，多面体却小了一圈。这类改写有个名字叫紧化，在实践中的收益经常大过换求解器或加机器。

再往上走会撞到天花板。近似难度（hardness of approximation）研究的是比例的下界：在 \(\text{P}\ne\text{NP}\) 之下，集合覆盖不可能在多项式时间内做到比 \(\ln n\) 更好；若唯一游戏猜想成立，最大割的 \(0.878\) 就是极限，Goemans--Williamson 已经到顶了。知道天花板在哪里，是为了不去花半年改进一个已经改无可改的东西。

至此，面对一个算不动的问题，手上的动作是成套的：先判断它落在难度地图的哪一格，再用松弛拿一个下界，然后在取整、随机舍入、分支定界之间挑一种把界换成方案，最后核对这个方案离最优有多远。这套动作里没有一步能凭空取消问题的困难。它们做的是另一件事——把``不知道离最优有多远''换成一个能写进交付说明的数字。
